\documentclass[]{deedy-resume-openfont}
\usepackage{fancyhdr}
 
\pagestyle{fancy}
\fancyhf{}
 
\begin{document}

%%%%%%%%%%%%%%%%%%%%%%%%%%%%%%%%%%%%%%
%
%     LAST UPDATED DATE
%
%%%%%%%%%%%%%%%%%%%%%%%%%%%%%%%%%%%%%%
%\lastupdated

%%%%%%%%%%%%%%%%%%%%%%%%%%%%%%%%%%%%%%
%
%     TITLE NAME
%
%%%%%%%%%%%%%%%%%%%%%%%%%%%%%%%%%%%%%%
\namesection{Matheus D.}{Krause}{ 
\href{mailto:krause.d.matheus@gmail.com}{krause.d.matheus@gmail.com} | +1 (515) 708-3842 | \href{https://mdkrause.github.io//}{about.me \ExternalLink} 
}

%%%%%%%%%%%%%%%%%%%%%%%%%%%%%%%%%%%%%%
%
%     COLUMN ONE
%
%%%%%%%%%%%%%%%%%%%%%%%%%%%%%%%%%%%%%%

\begin{minipage}[t]{1\textwidth} 

\sectionsep

Plant breeding is a fascinating and important science to ensure high quality and sufficient quantities of food are available for our growing population. This endeavor requires motivated scientists. In this scenario, my research interests comprise applied statistics, quantitative and statistical genetics. I am also excited about the future of agriculture, developing approaches to mitigate harmful environmental impact and improve efficiency at all levels of the value chain.

\sectionsep

%%%%%%%%%%%%%%%%%%%%%%%%%%%%%%%%%%%%%%
%     EDUCATION
%%%%%%%%%%%%%%%%%%%%%%%%%%%%%%%%%%%%%%

\section{Education} 

\sectionsep

\subsection{Ph.D. Student in Plant Breeding with Statistics minor (Exp. Summer 2023)}
\descript{Iowa State University (ISU) - Ames, Iowa, USA | Major Professor: Dr. William D. Beavis}
\project{Projects:} Development of
estimation methods to obtain unbiased realized genetic gain using data from routine field trials. \\ 
\hspace{1.55 cm} Development of novel breeding strategies to assure optimal trade-offs among
competing objectives.

\sectionsep
\vspace{0.2 mm}

\subsection{M.S in Plant Breeding and Genetics (2018)}
\descript{University of São Paulo (USP) - Piracicaba, SP, Brazil | Major Professor: Dr. Antonio A. F. Garcia}
\project{Title:} Boosting predictive ability of maize hybrids via genotype by environment interaction under multivariate GBLUP models.

\sectionsep
\vspace{0.2 mm}

\subsection{B.A in Agronomy (2015)}
\descript{Londrina State University (UEL) - Londrina, PR, Brazil | Major Professor: Dr. Josué M. Ferreira}
\project{Title:} Combining ability of maize inbred lines S$_9$ derived from ST15 and S709 synthetics.

\sectionsep
\vspace{0.2 mm}

\subsection{Exchange Student in Plant Breeding \& Engineering for Mediterranean and Tropical Areas (2013/14)}
\descript{Montpellier SupAgro - Montpellier, France | Under Drs. Ricardo Ralisch and Jean-Luc Regnard}
\project{Title:} Connaissance des Géniteurs Maïs (Confidential to Limagrain Europe and Montpellier SupAGro).

\sectionsep

%%%%%%%%%%%%%%%%%%%%%%%%%%%%%%%%%%%%%%
%     EXPERIENCE
%%%%%%%%%%%%%%%%%%%%%%%%%%%%%%%%%%%%%%

\section{Research Experience}

\sectionsep

\runsubsection{Bill Beavis's G.F. Sprague Quantitative Genetics Group}
\descript{| Graduate research assistant}
\project{May 2018 - Present | Department of Agronomy, ISU} \\
%\vspace{\topsep} 
- Develop a metric to estimate realized genetic gains on an annual basis using soybean routine field trials. \\
- Estimate density functions for variance components and partition of environments in homogeneous breeding regions.\\
- Strategies to assure optimal trade-offs among competing objectives for genetic 
improvement of soybean breeding programs. \\
- Develop software to correct partially informative markers from backcross one derived (BC1) double-haploid (DH) lines. \\
- Consult several graduate students (+15) with data analysis.

\sectionsep
\vspace{0.2 mm}

\runsubsection{Statistical Genetics Lab}
\descript{| Graduate research assistant}
\project{March 2016 - March 2018 | Department of Genetics, USP} \\
%\vspace{\topsep} 
- Apply multivariate linear mixed models to predict yield performance of single-cross maize hybrids based on information from genomics and genotype by environment interaction. 

\sectionsep
\vspace{0.2 mm}

\runsubsection{Tropical Breeding \& Genetics}
\descript{| Intern}
\project{Aug 2015 - Nov 2015 | Soybean Breeding} \\
%\vspace{\topsep}
- Plant phenotyping for soybean cyst nematode, \textit{Phytophthora sojae} and \textit{Sclerotinia sclerotiorum}, generation advancement in greenhouses, crossing blocks, seed packing, and general activities in the biotechnology laboratory.

\sectionsep
\vspace{0.2 mm}

\runsubsection{Laboratory of Maize Breeding}
\descript{| Undergraduate research assistant}
\project{Jun 2010 - Jun 2015 | Department of Biology, UEL} \\
%\vspace{\topsep} 
- Development and evaluation of inbred lines for common and sweet corn via classical and DH techniques. \\
- Evaluation and development of single-cross hybrids and synthetic populations for both common and sweet corn. \\
- Participation in an extension project to teach farmers how to produce their seeds of synthetic populations.

\sectionsep
\vspace{0.2 mm}

\runsubsection{Limagrain Europe}
\descript{| Intern}
\project{March 2014 - Jun 2014 | Department of Seed Production to Research} \\
%\vspace{\topsep}
- Data mining in a seed production database to support early decision-making in the maize breeding pipeline. \\
- Stability analysis of the checks used in routine breeder's seed field production. \\
- Adjustment in the methodology used to calculate future grain yield of the female inbred lines (i.g. its profitability for the company). 


\end{minipage} 
\hfill

\begin{minipage}[t]{1\textwidth} 

%%%%%%%%%%%%%%%%%%%%%%%%%%%%%%%%%%%%%%
%     PUBLICATIONS
%%%%%%%%%%%%%%%%%%%%%%%%%%%%%%%%%%%%%%

\section{Selected Journal Publications} 

\sectionsep
\textbf{Krause, MD}; Dias, KOG; Singh, AK; Beavis, WD. Using large soybean historical data to study genotype by environment variation and identify mega-environments with the integration of genetic and non-genetic factors. Submitted to Field Crop Research, 2022. \href{https://www.biorxiv.org/content/10.1101/2022.04.11.487885v3.abstract}{DOI: 10.1101/2022.04.11.487885}. \ExternalLink

\sectionsep
Dias, KOG; dos Santos, JPR; \textbf{Krause, MD}; Piepho, HP; Guimarães, LJM; Pastina, MM; Garcia, AAF. Leveraging probability concepts for cultivar recommendation in multi-environment trials. Theoretical and Applied Genetics, 2022. \href{https://link.springer.com/article/10.1007/s00122-022-04041-y}{DOI: 10.1007/s00122-022-04041-y}. \ExternalLink 

\sectionsep
Montes, CR; Fox, C; Sanz-Sáez, A; Serbin, SP; Kumagai, E; \textbf{Krause, MD}; Xavier, A; Specht, J; Beavis, WD; Bernacchi, CJ; Diers, BW; Ainsworth, EA. High-throughput characterization, correlation, and mapping of leaf photosynthetic and functional traits in the soybean (\textit{Glycine max}) nested association mapping population. Genetics, 2022. \href{https://academic.oup.com/genetics/advance-article/doi/10.1093/genetics/iyac065/6572345?login=true}{DOI: 10.1093/genetics/iyac065}. \ExternalLink

\sectionsep
Verzegnazzi, AL; Santos, I; \textbf{Krause, MD}; Hufford, M; Frei, UK; Campbell, J; Almeida, VC; Zuffo, LT; Boerman, N; Lübberstedt, T. Major locus for spontaneous haploid genome doubling detected by a case–control GWAS in exotic maize germplasm. Theoretical and Applied Genetics, 2021. \href{https://link.springer.com/article/10.1007/s00122-021-03780-8}{DOI: 10.1002/csc2.20253}. \ExternalLink 

\sectionsep
\textbf{Krause, MD}; Dias, KOG; dos Santos, JPR; Oliveira, AAO; Guimarães, LJM; Pastina, MM; Margarido, GRA; Garcia, AAF. Boosting predictive ability of tropical maize hybrids via genotype by environment interaction under multivariate GBLUP models. Crop Science, 2020. \href{https://acsess.onlinelibrary.wiley.com/doi/full/10.1002/csc2.20253}{DOI: 10.1002/csc2.20253}. \ExternalLink 

\sectionsep
Sekiya, A; Pestana, JK; da Silva, MGB; \textbf{Krause, MD}; da Silva, CRM; Ferreira, JM. Tropical supersweet corn haploid induction and ploidy determination at seedling stages. Brazilian Journal of Agricultural Research, 2020. \href{https://www.scielo.br/pdf/pab/v55/1678-3921-pab-55-e00968.pdf}{DOI: 10.1002/csc2.20253}. \ExternalLink 

\sectionsep
Xavier, LFS; Pestana, JK; Sekiya, A; \textbf{Krause, MD}; Moreira, RMP; Ferreira, JM. Partial diallel and potential of super sweet corn inbred lines bt$_2$ to obtain hybrids. Horticultura Brasileira, 2019. \href{http://www.scielo.br/scielo.php?script=sci_arttext&pid=S0102-05362019000300278}{DOI: 10.1590/s0102-053620190305}. \ExternalLink 

\sectionsep
Koltun, A; Cavalcante, AP; Lopes, KBA; \textbf{Krause, MD}; Marino, TP; Oliveira, ALM; Ferreira JM. Performance of maize hybrids from a partial diallel in association with \textit{Azospirillum}. African Journal of Agricultural Research, 2018. \href{https://academicjournals.org/journal/AJAR/article-abstract/B4A2A1B57541}{DOI: 10.5897/ajar2018.13077}. \ExternalLink

\sectionsep
\project{+ 25 Abstracts and 40 technical publications in agricultural newspapers.}

\sectionsep

%%%%%%%%%%%%%%%%%%%%%%%%%%%%%%%%%%%%%%
%     SKILLS
%%%%%%%%%%%%%%%%%%%%%%%%%%%%%%%%%%%%%%

\section{Additional information}

\sectionsep

\descript{Teaching \& Seminar}

\begin{tightemize}
	\vspace{3 mm}
	\item \project{2020/22} Lecture on \emph{AlphaSimR} in Agron 523 (Molecular P. Breed., graduate-level course, Dr. Lübberstedt). 
	\item \project{2021} Invited to speak in the ``Big Data: Manage your data before your data kills you" session at the Plant and Animal Genome Conference Plant and Animal Genome Conference (PAG) 2022 (declined due to COVID travel restrictions).
	\item \project{2017} Instructor of \emph{R programming} course at USP.
\end{tightemize}

\sectionsep

\descript{Computational Skills}
R programming language, basic Python and Julia, Linux, \LaTeX

\sectionsep

\descript{Language}
Portuguese (native), English, basic French

\sectionsep

\descript{Leadership}
\begin{tightemize}
	\item \project{2020-2021} Committee member in the 7th and 8th annual R. F. Baker Plant Breeding Symposium at ISU.
	\item \project{2011-2013} Student representative of Agronomy in the Regional Council of Engineering and Agronomy.
\end{tightemize}

\sectionsep

\descript{Awards}
C. R. Weber Award for Excellence in Plant Breeding (Spring 2022) \textbullet{} Bayer Travel Award (2020) \\

\sectionsep

\descript{Matheus enjoys...}
Spending time with his family and friends \textbullet{} Reading \textbullet{} Fishing \textbullet{} Being outdoors \textbullet{} Playing guitar \textbullet{} Riding motorcycle. \\

\sectionsep
\descript{For more information, please access}

\sectionsep
\sectionsep

% href with includegraphics works with lualatex

\hspace{6 cm} \href{https://www.linkedin.com/in/matheus-dalsente-krause-3012b158/}{\includegraphics[scale=0.125]{./logos/lin.jpeg}} \hspace{2 cm} \href{https://about.me/matheusdalsentekrause}{\includegraphics[scale=0.025]{./logos/g.png}} \hspace{2 cm} \href{https://mdkrause.github.io//}{\includegraphics[scale=0.06]{./logos/github.png}}

\end{minipage} 
\end{document}  \documentclass[]{article}
